\documentclass[11pt]{article}

\usepackage[margin=1in]{geometry}
\usepackage{graphicx}
\usepackage{array}
\usepackage{hyperref}
\usepackage{enumitem}
\usepackage{booktabs}
\usepackage{caption}
\usepackage{amsmath}
\usepackage{titlesec}

\titleformat{\section}{\large\bfseries}{\thesection.}{1em}{}

\begin{document}

\begin{center}
    \Large \textbf{Sri Sivasubramaniya Nadar College of Engineering, Chennai} \\
    \large (An Autonomous Institution affiliated to Anna University) \\
    \vspace{0.3cm}
\end{center}

\begin{table}[h!]
\renewcommand{\arraystretch}{1.5}
\centering
\resizebox{\textwidth}{!}{%
\begin{tabular}{|l|cll|}
\hline
\textbf{Degree \& Branch} & \multicolumn{1}{c|}{B.E Computer Science \& Engineering} & \textbf{Semester} & VI \\ \hline
\textbf{Subject Code \& Name} & \multicolumn{3}{c|}{UCS2612 -- Machine Learning Laboratory \& PREM SHANKAR.S} \\ \hline
\textbf{Academic Year} & \multicolumn{1}{c|}{2025--2026 (Even)} & \textbf{Batch} & 2023--2027 \\ \hline
\end{tabular}
}
\end{table}

\begin{center}
    \textbf{\Large Experiment 6}
\end{center}

\begin{center}
    \textbf{Decision Tree and Random Forest: A Comparative Classification Study}
\end{center}

%--------------------------------------------------

\section*{Objective}
\begin{itemize}
    \item To implement a \textbf{Decision Tree classifier}.
    \item To extend the Decision Tree into a \textbf{Random Forest ensemble model}.
    \item To study the impact of hyperparameters on overfitting and generalization.
    \item To \textbf{select optimal hyperparameters using 5-Fold Cross-Validation}.
    \item To compare single-tree and ensemble-tree models.
\end{itemize}

%--------------------------------------------------

\section*{Dataset}
\textbf{Wisconsin Diagnostic Breast Cancer Dataset}

\begin{itemize}
    \item Total samples: 569
    \item Features: 30 numerical attributes
    \item Target classes: Malignant (M) and Benign (B)
\end{itemize}

\textbf{Dataset Link:}  
\href{https://archive.ics.uci.edu/dataset/17/breast+cancer+wisconsin+diagnostic}{https://archive.ics.uci.edu}

%--------------------------------------------------

\section*{Theory}

\subsection*{Decision Tree Classifier}
A Decision Tree is a supervised learning model that recursively splits the feature space using impurity measures to form decision rules.

\textbf{Key Concepts:}
\begin{itemize}
    \item Gini Index and Entropy
    \item Tree depth and node splitting
    \item Overfitting in deep trees
\end{itemize}

\textbf{Limitation:}  
Decision Trees are prone to high variance and overfitting.

\subsection*{Random Forest Classifier}
Random Forest is an ensemble learning technique that combines multiple decision trees trained on bootstrapped samples.

\textbf{Key Ideas:}
\begin{itemize}
    \item Bootstrap aggregation (bagging)
    \item Random feature selection
    \item Majority voting
\end{itemize}

\textbf{Advantage:}  
Reduces variance and improves generalization compared to a single Decision Tree.

%--------------------------------------------------

\section*{Steps for Implementation}

\begin{enumerate}[label=\arabic*.]
    \item Load the dataset and encode class labels.
    \item Perform Exploratory Data Analysis:
    \begin{itemize}
        \item Class distribution
        \item Feature correlation analysis
    \end{itemize}
    \item Split the dataset into training and testing sets (80--20).
    \item Train a Decision Tree classifier.
    \item Define a hyperparameter search space.
    \item Perform \textbf{5-Fold Cross-Validation} to evaluate each hyperparameter combination.
    \item Select the hyperparameters that yield the best average cross-validation performance.
    \item Train a Random Forest classifier.
    \item Repeat cross-validation-based hyperparameter selection.
    \item Compare both models using evaluation metrics.
\end{enumerate}

%--------------------------------------------------

\section*{Hyperparameters to be Explored}

\subsection*{Decision Tree}
\begin{itemize}
    \item \texttt{criterion}
    \item \texttt{max\_depth}
    \item \texttt{min\_samples\_split}
    \item \texttt{min\_samples\_leaf}
\end{itemize}

\subsection*{Random Forest}
\begin{itemize}
    \item \texttt{n\_estimators}
    \item \texttt{max\_depth}
    \item \texttt{max\_features}
    \item \texttt{bootstrap}
\end{itemize}

%--------------------------------------------------

\section*{Hyperparameter Tuning Results}

\subsection*{Decision Tree Cross-Fold Results}
\begin{table}[h!]
\centering
\caption{Decision Tree Hyperparameter Evaluation using 5-Fold Cross-Validation}
\begin{tabular}{cccc}
\toprule
Criterion & Max Depth & Avg CV Accuracy (\%) & Avg CV F1 Score \\
\midrule
 &  &  &  \\
 &  &  &  \\
 &  &  &  \\
\bottomrule
\end{tabular}
\end{table}

\subsection*{Random Forest Cross-Fold Results}
\begin{table}[h!]
\centering
\caption{Random Forest Hyperparameter Evaluation using 5-Fold Cross-Validation}
\begin{tabular}{ccccc}
\toprule
n\_estimators & Max Depth & Max Features & Avg CV Accuracy (\%) & Avg CV F1 Score \\
\midrule
 &  &  &  &  \\
 &  &  &  &  \\
 &  &  &  &  \\
\bottomrule
\end{tabular}
\end{table}

%--------------------------------------------------

\section*{5-Fold Cross-Validation Performance Comparison}

\begin{table}[h!]
\centering
\caption{5-Fold Cross-Validation Accuracy Comparison}
\begin{tabular}{lcccccc}
\toprule
Model & Fold 1 & Fold 2 & Fold 3 & Fold 4 & Fold 5 & Average \\
\midrule
Decision Tree &  &  &  &  &  &  \\
Random Forest &  &  &  &  &  &  \\
\bottomrule
\end{tabular}
\end{table}

%--------------------------------------------------

\section*{Evaluation Metrics}
\begin{itemize}
    \item Accuracy
    \item Precision
    \item Recall
    \item F1-score
    \item Confusion Matrix
    \item ROC Curve and AUC
\end{itemize}

%--------------------------------------------------

\section*{Observation Questions}
\begin{itemize}
    \item How does tree depth affect overfitting in Decision Trees?
    \item Which hyperparameter had the greatest impact on performance?
    \item How does Random Forest improve generalization?
    \item Did ensemble learning always improve performance? Why or why not?
\end{itemize}

%--------------------------------------------------

\section*{Conclusion}
Decision Tree and Random Forest models were implemented and evaluated using 5-fold cross-validation.  
Hyperparameters were selected based on average cross-validation performance, ensuring robust generalization.  
The results demonstrate that Random Forest reduces variance and improves stability compared to a single Decision Tree.

%--------------------------------------------------

\section*{References}
\begin{itemize}
    \item \href{https://scikit-learn.org/stable/modules/tree.html}{Scikit-learn: Decision Trees}
    \item \href{https://scikit-learn.org/stable/modules/ensemble.html}{Scikit-learn: Random Forest}
    \item \href{https://archive.ics.uci.edu/dataset/17/breast+cancer+wisconsin+diagnostic}{UCI Dataset}
\end{itemize}

\end{document}
